\section{Motivation}

The determinant compresses an entire $n\times n$ matrix into a single scalar that answers a remarkable
number of yes/no questions at once. Is the matrix invertible? Does the system $AX=B$ have exactly one
solution? Do these $n$ vectors form a basis of $\K^n$? Does the transformation collapse space onto
something flatter? All of these questions are equivalent, and all reduce to checking $\det(A)\neq 0$ —
which is why a determinant test appears in virtually every algorithm that must know whether a matrix is
``degenerate''.

For a computer scientist the determinant has a second life as a \emph{geometric} quantity:

\begin{itemize}
\item In 2D, $|\det(A)|$ is the \textbf{area} of the parallelogram spanned by the rows of $A$; in 3D it is
a volume. A transformation with determinant $6$ magnifies every area six-fold; determinant $0$ means space
gets flattened.
\item The \textbf{sign} of the determinant detects \emph{orientation}. The test ``does point $C$ lie to
the left or to the right of the line through $A$ and $B$?'' is the sign of a $2\times 2$ determinant; it
is the basic primitive of computational geometry, executed millions of times inside convex-hull
algorithms, triangulations, collision detection and map matching. In 3D rendering the same sign test
(\emph{backface culling}) lets the GPU skip the triangles that face away from the camera.
\item In graphics and machine learning, determinants of \emph{Jacobian} matrices track how transformations
distort volume — the bookkeeping device behind texture mapping and behind generative ``normalizing flow''
models.
\end{itemize}

Part of the motivation is also a computational warning. The Laplace expansion defined below is
mathematically fundamental but computationally explosive: applied naively it performs on the order of $n!$
operations. The behaviour of determinants under row operations, established in this chapter, brings the
cost down to that of Gaussian elimination, $O(n^3)$ — one of the nicest examples in this course of theory
rescuing an algorithm. The same warning applies to the chapter's elegant closed formulas (Cramer's rule
and the adjugate formula for $A^{-1}$): they are invaluable for insight and perfectly practical for small
matrices, but for large ones the method of choice is, once again, row reduction.

\section{Determinants}

To every square matrix $A \in \K^n_n$ over a field $\K$ one assigns a scalar from $\K$, denoted $|A|$ or
$\det(A)$ and called \emph{the determinant of $A$}. For matrices of sizes $1\times 1$, $2\times 2$ and
$3\times 3$ it is given by the explicit formulas
\[
\begin{aligned}
| a_{11} | &= a_{11},\\[2pt]
\left | \begin{array}{cc} a_{11} & a_{12}\\ a_{21} & a_{22} \end{array} \right | &= a_{11} a_{22} - a_{12} a_{21},\\[2pt]
\left | \begin{array}{ccc}
a_{11} & a_{12} & a_{13}\\ a_{21} & a_{22} & a_{23}\\ a_{31} & a_{32} & a_{33}
\end{array} \right | &=
a_{11} a_{22} a_{33} + a_{12} a_{23} a_{31}+a_{13} a_{21} a_{32}
- a_{13}a_{22}a_{31} -a_{12} a_{21} a_{33} - a_{11} a_{23} a_{32}.
\end{aligned}
\]

\begin{example}
\[
\left| \begin{array}{cc} 1 & 2\\ 3 & 4 \end{array} \right| = 1\cdot 4 - 2\cdot 3 = -2,
\]
and, applying the $3\times 3$ formula term by term,
\[
\left| \begin{array}{ccc} 2 & 1 & 3\\ 0 & 1 & -1\\ 1 & 2 & 1 \end{array} \right|
= 2\cdot 1\cdot 1 + 1\cdot(-1)\cdot 1 + 3\cdot 0\cdot 2
- 3\cdot 1\cdot 1 - 1\cdot 0\cdot 1 - 2\cdot(-1)\cdot 2
= 2 - 1 + 0 - 3 - 0 + 4 = 2.
\]
\end{example}

\begin{definition}[Laplace expansion]
Let $A$ be an $n\times n$ matrix and let $A_{ij}$ denote the matrix obtained from $A$ by deleting the $i$-th
row and the $j$-th column. For any fixed column index $j$ between $1$ and $n$,
\[
\det(A) = \sum_{i=1}^{n} (-1)^{i+j} a_{ij}\cdot \det(A_{ij}).
\]
This formula is the \emph{Laplace expansion along the $j$-th column}. In particular, for $j=1$,
\[
\det(A) = \sum_{i=1}^{n} (-1)^{i+1} a_{i 1}\cdot \det(A_{i 1}).
\]
\end{definition}

\begin{remark}
Since $\det(A) = \det(A^T)$ (see the properties below), the determinant may equally well be expanded
along any fixed \emph{row}: $\det(A) = \sum_{j=1}^{n} (-1)^{i+j} a_{ij}\cdot\det(A_{ij})$ for any fixed
row index $i$. In practice one expands along the row or column containing the most zeros.
\end{remark}

\begin{example}
We compute the determinant by expanding along the first column:
\[
\left | \begin{array}{ccc}
1 & 2 & 3\\ 4 & 5 & 6\\ 7 & 8 & 9
\end{array} \right |
= 1\cdot\left | \begin{array}{cc} 5 & 6\\ 8 & 9 \end{array} \right |
- 4 \cdot \left | \begin{array}{cc} 2 & 3\\ 8 & 9 \end{array} \right |
+ 7\cdot \left | \begin{array}{cc} 2 & 3\\ 5 & 6 \end{array} \right | = 0.
\]
\end{example}

\begin{fact}[Properties of determinants]
Let $A$ be a square matrix of size $n\times n$. Then
\begin{enumerate}
\item $\det(A) = \det(A^T)$;
\item $\det(A) = 0$ if $A$ has a zero row (or column) or two identical rows (or columns);
\item $\det(A) = 0$ if and only if $\rank(A)<n$;
\item $\det(A) = -\det(B)$ if $B$ is obtained from $A$ by a single row switch ($R_i\leftrightarrow R_j$);
\item $\det(A) = \det(B)$ if $B$ is obtained from $A$ by a single row addition
($R_i+k\cdot R_j\rightarrow R_i$);
\item $\det(B) = k\cdot \det(A)$ if $B$ is obtained from $A$ by the row scaling $kR_i\rightarrow R_i$.
\end{enumerate}
The last three properties also hold for the corresponding elementary column operations.
\end{fact}

\begin{example}
Using row operations to reach a triangular form:
\[
\left | \begin{array}{ccc}
1 & 2 & 3\\ 4 & 5 & 6\\ 7 & 8 & 9
\end{array} \right |
\overset{R_2-4 R_1}{=}
\left | \begin{array}{ccc}
1 & 2 & 3\\ 0 & -3 & -6\\ 7 & 8 & 9
\end{array} \right |
\overset{R_3-7 R_1}{=}
\left | \begin{array}{ccc}
1 & 2 & 3\\ 0 & -3 & -6\\ 0 & -6 & -12
\end{array} \right |
= 2\cdot \left | \begin{array}{ccc}
1 & 2 & 3\\ 0 & -3 & -6\\ 0 & -3 & -6
\end{array} \right | = 2\cdot 0 = 0.
\]
\end{example}

\begin{fact}
The determinant of an upper (or lower) triangular matrix equals the product of the entries on its main
diagonal. In particular this applies to diagonal matrices, and $\det( I_{n\times n} ) = 1$.
\end{fact}

\begin{example}
\[
\left| \begin{array}{ccc} 2 & 7 & -1\\ 0 & 3 & 5\\ 0 & 0 & -4 \end{array} \right| = 2\cdot 3\cdot (-4) = -24.
\]
Combined with row reduction, this fact gives an efficient general method of computing determinants: reduce
the matrix to a triangular form, keeping track of how each operation changes the determinant, exactly as in
the example above.
\end{example}

\begin{example}[A $4\times 4$ determinant]
Row additions do not change the determinant, so
\[
\left| \begin{array}{cccc}
1 & 2 & 0 & 1\\ 2 & 5 & 2 & 3\\ 1 & 2 & 1 & 2\\ 3 & 6 & 0 & 4
\end{array} \right|
\overset{\substack{R_2-2R_1\\ R_3-R_1\\ R_4-3R_1}}{=}
\left| \begin{array}{cccc}
1 & 2 & 0 & 1\\ 0 & 1 & 2 & 1\\ 0 & 0 & 1 & 1\\ 0 & 0 & 0 & 1
\end{array} \right| = 1\cdot 1\cdot 1\cdot 1 = 1.
\]
For larger matrices this approach is dramatically faster than a direct Laplace expansion.
\end{example}

\begin{fact}[Multiplicativity]
Let $A$ and $B$ be square matrices of size $n\times n$. Then $\det(A\cdot B) = \det(A)\cdot \det(B)$.
\end{fact}

\begin{warning}
The determinant is \emph{not} additive: in general $\det(A+B)\neq \det(A)+\det(B)$. Note also that
$\det(k\cdot A) = k^n\det(A)$ for an $n\times n$ matrix $A$: scaling the whole matrix scales \emph{each} of
its $n$ rows.
\end{warning}

\begin{remark}[Geometric meaning]
For a real $2\times 2$ matrix $A$, the absolute value $|\det(A)|$ is the area of the parallelogram spanned
by the two rows (equivalently, the two columns) of $A$; the sign of the determinant records the orientation
of this pair of vectors. Similarly, for a real $3\times 3$ matrix, $|\det(A)|$ is the volume of the
parallelepiped spanned by the rows.
\end{remark}

\section{Determinants and systems of linear equations}

\begin{theorem}
Let $AX=B$ be a system of $n$ linear equations with $n$ unknowns. Then $AX=B$ has a unique solution if and
only if $\det(A)\neq 0$.
\end{theorem}

\begin{theorem}[Cramer's rule]
Consider a system $AX=B$ with $n$ equations and $n$ unknowns, where $A$ is a square matrix of size
$n\times n$. Let $A_{|i}$ denote the matrix obtained from $A$ by replacing its $i$-th column with the column
$B$. If $\det(A)\neq 0$, then
\[
x_1 = \frac{\det(A_{|1})}{\det(A)}, \qquad
x_2 = \frac{\det(A_{|2})}{\det(A)}, \qquad
\ldots, \qquad
x_n = \frac{\det(A_{|n})}{\det(A)}.
\]
\end{theorem}

\begin{example}
Consider the system $\{\,x+y = 1,\ x+2y = 0\,\}$. Here
\[
A = \left ( \begin{array}{cc} 1 & 1\\ 1 & 2 \end{array} \right ), \quad \det(A) = 1,
\]
\[
A_{|1} = \left ( \begin{array}{cc} 1 & 1\\ 0 & 2 \end{array} \right ), \ \det(A_{|1}) = 2, \qquad
A_{|2} = \left ( \begin{array}{cc} 1 & 1\\ 1 & 0 \end{array} \right ), \ \det(A_{|2}) = -1.
\]
Hence $x=\tfrac{2}{1}=2$ and $y=\tfrac{-1}{1}=-1$.
\end{example}

\begin{example}[Cramer's rule for a $3\times 3$ system]
Consider the system
\[
\begin{aligned}
x+y+z &= 2,\\ x-y+z &= 0,\\ x+y-z &= 0.
\end{aligned}
\]
We compute
\[
\det(A) = \left| \begin{array}{ccc} 1 & 1 & 1\\ 1 & -1 & 1\\ 1 & 1 & -1 \end{array} \right| = 4,
\qquad
\det(A_{|1}) = \left| \begin{array}{ccc} 2 & 1 & 1\\ 0 & -1 & 1\\ 0 & 1 & -1 \end{array} \right| = 0,
\]
\[
\det(A_{|2}) = \left| \begin{array}{ccc} 1 & 2 & 1\\ 1 & 0 & 1\\ 1 & 0 & -1 \end{array} \right| = 4,
\qquad
\det(A_{|3}) = \left| \begin{array}{ccc} 1 & 1 & 2\\ 1 & -1 & 0\\ 1 & 1 & 0 \end{array} \right| = 4.
\]
Hence $x = \tfrac04 = 0$, $y = \tfrac44 = 1$ and $z = \tfrac44 = 1$; a direct check confirms that $(0,1,1)$
solves the system.
\end{example}

\section{Matrix inversion}

\begin{definition}
Let $A$ be a square matrix of size $n\times n$. A matrix $B$ of the same size is called an \emph{inverse} of
$A$ if
\[
A\cdot B = I_{n\times n}, \qquad B\cdot A = I_{n\times n}.
\]
If $A$ has an inverse, it is unique; it is denoted by $A^{-1}$.
\end{definition}

\begin{example}
For $A = \left ( \begin{smallmatrix} 1 & 0\\ 0 & 2 \end{smallmatrix} \right )$ one checks that
$A^{-1} = \left ( \begin{smallmatrix} 1 & 0\\ 0 & \tfrac12 \end{smallmatrix} \right )$.
\end{example}

\begin{warning}
Not every matrix is invertible.
\end{warning}

\begin{theorem}
A square matrix $A$ is invertible if and only if $\det(A)\neq 0$. If $A$ is invertible, then
\[
\det(A^{-1}) = \frac{1}{\det(A)}.
\]
\end{theorem}

\begin{proof}
We prove the second statement. Since $\det(I) = 1$ and $\det(A\cdot B) = \det(A)\cdot \det(B)$, and since
$A\cdot A^{-1} = I$ for invertible $A$, we get
\[
1 = \det(I) = \det(A\cdot A^{-1}) = \det(A)\cdot \det(A^{-1}). \qedhere
\]
\end{proof}

\begin{theorem}[Properties of the inverse]
Let $A,B$ be invertible matrices of size $n\times n$ and let $k\neq 0$ be a scalar. Then $A^{-1}$,
$A\cdot B$, $A^T$ and $k\cdot A$ are all invertible, and
\[
(A^{-1})^{-1} = A, \qquad (A\cdot B)^{-1} = B^{-1}\cdot A^{-1}, \qquad (A^T)^{-1} = (A^{-1})^T,
\qquad (k\cdot A)^{-1} = \tfrac1k\cdot A^{-1}.
\]
\end{theorem}

\begin{warning}
Note the reversed order in $(A\cdot B)^{-1} = B^{-1}\cdot A^{-1}$ — the same reversal as in the transpose
formula $(A\cdot B)^T = B^T\cdot A^T$.
\end{warning}

\begin{fact}[Inversion by row reduction]
Let $A$ be a square matrix of size $n\times n$. Form the $n\times 2n$ matrix $C=(A\mid I_{n\times n})$ and row reduce
it to the form $(I_{n\times n}\mid B)$. Such a reduction is possible if and only if $A$ is invertible, and
the resulting matrix $B$ is the inverse of $A$.
\end{fact}

\begin{example}
\[
\left ( \begin{array}{cc|cc}
1 & 0 & 1 & 0 \\ 2 & 2 & 0 & 1
\end{array} \right )
\xrightarrow{R_2-2R_1\to R_2}
\left ( \begin{array}{cc|cc}
1 & 0 & 1 & 0 \\ 0 & 2 & -2 & 1
\end{array} \right )
\xrightarrow{\frac{1}{2}R_2\to R_2}
\left ( \begin{array}{cc|cc}
1 & 0 & 1 & 0 \\ 0 & 1 & -1 & \tfrac{1}{2}
\end{array} \right ),
\]
so $\left ( \begin{smallmatrix} 1 & 0\\ 2 & 2 \end{smallmatrix} \right )^{-1}
= \left ( \begin{smallmatrix} 1 & 0\\ -1 & \tfrac{1}{2} \end{smallmatrix} \right )$.
\end{example}

\begin{example}[A $3\times 3$ inverse]
We invert $A = \left ( \begin{smallmatrix} 1 & 0 & 1\\ 1 & 1 & 0\\ 0 & 1 & 1 \end{smallmatrix} \right )$:
\[
\left ( \begin{array}{ccc|ccc}
1 & 0 & 1 & 1 & 0 & 0\\ 1 & 1 & 0 & 0 & 1 & 0\\ 0 & 1 & 1 & 0 & 0 & 1
\end{array} \right )
\xrightarrow{R_2-R_1\to R_2}
\left ( \begin{array}{ccc|ccc}
1 & 0 & 1 & 1 & 0 & 0\\ 0 & 1 & -1 & -1 & 1 & 0\\ 0 & 1 & 1 & 0 & 0 & 1
\end{array} \right )
\xrightarrow{R_3-R_2\to R_3}
\left ( \begin{array}{ccc|ccc}
1 & 0 & 1 & 1 & 0 & 0\\ 0 & 1 & -1 & -1 & 1 & 0\\ 0 & 0 & 2 & 1 & -1 & 1
\end{array} \right )
\]
\[
\xrightarrow{\frac12 R_3\to R_3}
\left ( \begin{array}{ccc|ccc}
1 & 0 & 1 & 1 & 0 & 0\\ 0 & 1 & -1 & -1 & 1 & 0\\ 0 & 0 & 1 & \tfrac12 & -\tfrac12 & \tfrac12
\end{array} \right )
\xrightarrow{\substack{R_2+R_3\to R_2\\ R_1-R_3\to R_1}}
\left ( \begin{array}{ccc|ccc}
1 & 0 & 0 & \tfrac12 & \tfrac12 & -\tfrac12\\
0 & 1 & 0 & -\tfrac12 & \tfrac12 & \tfrac12\\
0 & 0 & 1 & \tfrac12 & -\tfrac12 & \tfrac12
\end{array} \right ),
\]
so
\[
A^{-1} = \frac12 \left ( \begin{array}{ccc} 1 & 1 & -1\\ -1 & 1 & 1\\ 1 & -1 & 1 \end{array} \right ).
\]
One can (and always should) verify the result by checking that $A\cdot A^{-1} = I$.
\end{example}

\begin{fact}[Inversion by determinants]
Let $A$ be an invertible square matrix and let $A_{k,m}$ denote the matrix obtained from $A$ by deleting its
$k$-th row and $m$-th column. Then
\[
A^{-1} = \frac{1}{|A|}\cdot
\left( \begin{array}{cccc}
(-1)^{1+1} |A_{11}| & (-1)^{1+2} |A_{12}| & \ldots & (-1)^{1+n} |A_{1n}| \\
(-1)^{2+1} |A_{21}| & (-1)^{2+2} |A_{22}| & \ldots & (-1)^{2+n} |A_{2n}| \\
\vdots & \vdots & \ddots & \vdots \\
(-1)^{n+1} |A_{n1}| & (-1)^{n+2} |A_{n2}| & \ldots & (-1)^{n+n}|A_{nn}|
\end{array} \right)^T.
\]
\end{fact}

\begin{example}
\[
\left ( \begin{array}{cc} 1 & 0 \\ 2 & 2 \end{array} \right )^{-1}
= \frac{1}{2} \left ( \begin{array}{cc}
|(2)| & -|(2)| \\ -|(0)| & |(1)|
\end{array} \right )^T
= \frac{1}{2} \left ( \begin{array}{cc} 2 & 0 \\ -2 & 1 \end{array} \right )
= \left ( \begin{array}{cc} 1 & 0 \\ -1 & \tfrac12 \end{array} \right ).
\]
\end{example}

\begin{fact}[Inverse of a $2\times 2$ matrix]
Let $A= \left ( \begin{smallmatrix} a & b \\ c & d \end{smallmatrix} \right )$ with
$\det(A) = ad-bc \neq 0$. Then
\[
A^{-1} = \frac{1}{ad-bc}\cdot \left( \begin{array}{cc} d & -b \\ -c & a \end{array} \right ).
\]
\end{fact}
